% OakGPT chat evaluation v1 report. LaTeX source of report.md; the two must stay in sync.
% Build (pdfLaTeX, XeLaTeX or Tectonic) from this directory, e.g. `tectonic report.tex`.
% The source is pure ASCII so every engine accepts it.
\documentclass[11pt,a4paper]{article}

\usepackage{iftex}
\ifPDFTeX
  \usepackage[T1]{fontenc}
  \usepackage{lmodern}
\fi
\usepackage[margin=25mm]{geometry}
\usepackage{microtype}
\usepackage{booktabs}
\usepackage{tabularx}
\usepackage{array}
\usepackage{xcolor}
\usepackage{enumitem}
\usepackage{fancyvrb}
\usepackage{caption}
\usepackage{float}
\usepackage[hidelinks]{hyperref}

\definecolor{oak}{HTML}{3D6B35}
\hypersetup{colorlinks=true,linkcolor=oak,urlcolor=oak,citecolor=oak,
  pdftitle={OakGPT Chat Evaluation v1: How Good Are the Local Models?},
  pdfauthor={OakGPT - Local LLM}}
\setlist{itemsep=2pt,topsep=4pt}
\setlength{\parindent}{0pt}
\setlength{\parskip}{6pt}
\newcolumntype{Y}{>{\raggedright\arraybackslash}X}
\newcolumntype{R}{>{\raggedleft\arraybackslash}X}
\renewcommand{\arraystretch}{1.15}
\urlstyle{same}
\setlength{\emergencystretch}{3em}

\newcommand{\code}[1]{\texttt{#1}}
\newcommand{\decision}[1]{\subsection*{#1}\addcontentsline{toc}{subsection}{#1}}
\newcommand{\field}[1]{\textbf{#1}\quad}

\title{\textbf{OakGPT Chat Evaluation v1:\\How Good Are the Local Models?}\\[4pt]
  \large Design decisions, rationale, sources and measured results}
\author{OakGPT -- Local LLM}
\date{2026-11-23}

\begin{document}
\maketitle

\begin{center}
\small
\begin{tabularx}{0.92\textwidth}{@{}lY@{}}
\toprule
Track & Chat, evaluation before C2 (fine-tuning) \\
Branch & \path{feature/chat-c1-local-web-chat} \\
Dataset & \path{data/eval/chat_eval_v1}, frozen in commit \code{851c057};
  \code{items.jsonl} SHA-256 \code{61946eb4...2ecb} \\
Raw results & \path{chat_eval_v1_results.json}: 441 replies with scores, timings, model
  hashes and environment \\
Hardware & NVIDIA GeForce RTX 2060, 6\,GB VRAM, Windows 10 19045 \\
Inference & llama.cpp \code{b11930-85cc432e2}; Qwen3 1.7B Q4\_K\_M (text),
  Qwen3-VL-2B-Instruct Q4\_K\_M + F16 encoder (images) \\
\bottomrule
\end{tabularx}
\end{center}

\begin{abstract}
\noindent Until now, OakGPT's answer quality was only described by anecdotes. This report
adds a frozen evaluation set of 67 items (40 text, 27 image) with answer keys, a runner
that uses the app's own code path, and one complete measurement: 441 replies, three fixed
seeds per item and condition. \emph{Strict} means the reply equals the key in the
requested form.

Thinking mode raises the text model's strict accuracy from 49.2\% to 91.7\% (paired
difference 42.5 points, 95\% interval 27.5--56.7). The cost is a median of 1.9\,s
instead of 0.06\,s per short answer. In direct mode, the 2B vision model answers the
text items 21.7 points better than the 1.7B text model (9.2--34.2). The vision model was
always right on colour, shape, position and size, but read only 20\% of the pixel-font
text correctly. These numbers are the baseline that any C2 fine-tuning must beat; they
describe this small set only.
\end{abstract}

\begin{table}[H]
\centering\small
\begin{tabularx}{\textwidth}{@{}Yrllr@{}}
\toprule
Condition & Items & Strict accuracy & Lenient accuracy & Median time \\
\midrule
Text model, chat preset (UI default) & 40 & \textbf{49.2\%} [35.0, 64.2] & 60.8\% [45.8, 75.8] & 0.06\,s \\
Text model, reasoning preset & 40 & \textbf{91.7\%} [83.3, 98.3] & 91.7\% [83.3, 98.3] & 1.92\,s \\
Vision model, image items & 27 & \textbf{75.3\%} [60.5, 88.9] & 75.3\% [60.5, 88.9] & 0.06\,s \\
Vision model, text items & 40 & \textbf{70.8\%} [56.7, 84.2] & 75.8\% [63.3, 88.3] & 0.05\,s \\
\bottomrule
\end{tabularx}
\caption*{Item-level 95\% bootstrap intervals in brackets. Time is per reply.}
\end{table}

\tableofcontents

\section{What was built}

\begin{center}\small
\begin{tabularx}{\textwidth}{@{}Yl@{}}
\toprule
Part & File \\
\midrule
Frozen set (items, 29 PNGs, data card) & \path{data/eval/chat_eval_v1/} \\
Deterministic builder (standard library only) & \path{scripts/build_chat_eval_v1.py} \\
Runner, scoring and statistics & \path{src/oakgpt/chat/evaluate.py} \\
Command & \code{oakgpt chat-eval --config ... --dataset ... --output ...} \\
CPU tests (25, no model downloads) & \path{tests/test_chat_eval.py} \\
\bottomrule
\end{tabularx}
\end{center}

Text categories, with items per category: arithmetic 6, word problems 6, facts 6,
Swedish 6, reading Python code 7, explicit instructions 5, reading comprehension 4. Image
categories: counting 5, colour 4, shape 3, position 4, size 2, bar chart 2, OCR 5 and
two-image questions 2. Every prompt asks for a short answer in a stated form, for example
``Reply with only the number.''

\section{Decision records}

\decision{D25. Write a small project-specific set instead of running a public benchmark}
\field{Decision.} Write 67 items covering what OakGPT is used for: short factual,
arithmetic, code and Swedish questions, and basic image questions. Do not run MMLU- or
GSM8K-style suites as the first evaluation.

\field{Rationale.}
\begin{itemize}
\item Swedish and image questions are product features that common English text suites
  do not cover.
\item The whole set needs about 6 minutes of generation on the RTX 2060, so it can run
  after every change, including every C2 training run.
\item Public suites are widely published, so pretrained models may have seen them.
  Lessons on reproducible evaluation~\cite{biderman} show that scores also move with
  prompt format and answer extraction. That makes a borrowed number hard to compare with
  our own setup.
\end{itemize}

\field{Trade-off.} The set is small and written by one author, so it cannot rank OakGPT
against published models. A public suite can be added later as a separate, complementary
measurement.

\decision{D26. Synthetic images whose answers follow from their construction}
\field{Decision.} Draw every image from an explicit specification with the Python
standard library (a \code{zlib} PNG encoder and a built-in 5$\times$7 pixel font).
Examples: ``7 orange circles at seeded non-overlapping positions'' or ``bars A--D with
heights 120, 200, 310, 160''.

\field{Rationale.}
\begin{itemize}
\item The key is a property of the specification, not a person's reading of a photo, so
  it cannot be wrong through labelling mistakes.
\item No external images are needed, so there are no licence or privacy questions.
\item The builder is byte-for-byte deterministic. A test rebuilds the set and requires
  identical files, so any accidental change to the frozen data fails the test suite.
\end{itemize}

\field{Trade-off.} Flat shapes and a blocky pixel font are not photographs or ordinary
typefaces. Section~\ref{sec:interpretation} shows that the font matters.

\decision{D27. Freeze in Git before the first model run}
\field{Decision.} Commit the set, the builder and the runner (\code{851c057})
\emph{before} any model sees an item. Never change prompts, images or keys afterwards;
changes need \code{chat\_eval\_v2}.

\field{Rationale.} If keys or prompts were adjusted after reading model output, the set
would measure the author's adjustments rather than the models. The commit order in Git is
the evidence that this did not happen. The result file records the commit and the dataset
hash. It also records the dirty flag; here \code{git\_dirty} is true only because of four
unrelated untracked PDFs in \path{output/pdf/}.

\decision{D28. Rule-based strict and lenient scoring, no model as judge}
\field{Decision.} Score every reply in two ways. \emph{Strict}: NFKC normalisation, case
folding, whitespace collapse, and stripping of surrounding punctuation, quotes and
Markdown emphasis; the reply must then equal an accepted answer. \emph{Lenient}: an
accepted answer appears as a whole word anywhere in the reply. Only the answer channel
is scored, never thinking text.

\field{Rationale.}
\begin{itemize}
\item Normalising before exact match follows SQuAD's evaluation~\cite{squad}.
\item Strict scoring rewards what a user of a chat app wants: the right answer in the
  form they asked for.
\item Lenient scoring separates ``wrong'' from ``right but wordy''. For example, reading
  comprehension scores 25\% strict and 100\% lenient in chat mode.
\item An LLM judge was rejected for three reasons. Judges show position, verbosity and
  self-enhancement biases~\cite{judge}. A cloud judge would break the rule that prompts
  never leave the machine. A local 2B judge would be weaker than the models it grades.
\end{itemize}

\field{Trade-off.} Lenient scoring can pass hedged or self-contradicting replies, so all
raw replies are kept. Strict scoring can fail a correct answer in a reasonable but
unexpected form: \code{3,1} for the expected output \code{3 1} fails.

\decision{D29. Measure the product: the app's code path, presets and three fixed seeds}
\field{Decision.} Send every item as a fresh one-turn conversation, with the configured
system prompt, through the same \code{ModelPool}, templates and presets as the web UI.
Run each item with seeds 2026, 2027 and 2028, the same seeds in every condition. Do not
use greedy decoding.

\field{Rationale.}
\begin{itemize}
\item The point is to know what a user of OakGPT gets, so the UI's sampling settings are
  used rather than special evaluation settings.
\item The Qwen3 model card advises against greedy decoding in thinking mode because it
  can cause degradation and endless repetition~\cite{qwen3card}.
\item Fixed seeds make the run repeatable and pair the conditions item by item and seed
  by seed.
\end{itemize}

\field{Trade-off.} Sampling makes some items flip between seeds. In chat mode, 3 of 40
items had mixed results; three seeds are too few to estimate per-item reliability.

\decision{D30. Item-level bootstrap intervals, with repeats averaged first}
\field{Decision.} Average the three replies per item, then resample \emph{items} 2000
times with seed 2026 for a 95\% percentile interval. Condition differences are paired by
item.

\field{Rationale.} Three samples of the same prompt are not three independent questions;
treating them as such would shrink the intervals falsely. This follows the existing
OakGPT convention and the bootstrap method~\cite{efron}. Miller~\cite{miller} shows that
related questions, such as several questions on one passage, make naive intervals too
narrow. Here only two passages carry two questions each, and the shape items share
templates, so the reported intervals are somewhat optimistic.

\decision{D31. Also run the vision model on the text items}
\field{Decision.} Add a \code{vision-on-text} condition: the same 40 text items, answered
by Qwen3-VL-2B with the vision preset.

\field{Rationale.} In C1.1 (architecture overview and D16), any conversation with an
image stays on the vision model, including later text-only turns. That was a speed
decision whose effect on text quality was unknown. This condition measures it directly.

\section{Results}\label{sec:results}

Per-category strict accuracy, with lenient accuracy in brackets where it differs. Each
category has 2--7 items, so a difference of one item is 14--50 points.

\begin{table}[H]
\centering\small
\begin{tabularx}{\textwidth}{@{}YrRRR@{}}
\toprule
Text category & Items & Chat preset & Reasoning preset & Vision model \\
\midrule
Arithmetic & 6 & 83\% & 100\% & 83\% (100\%) \\
Word problems & 6 & 50\% & 89\% & 50\% (56\%) \\
Facts & 6 & 83\% & 100\% & 100\% \\
Swedish & 6 & 6\% (33\%) & 83\% & 44\% \\
Python code & 7 & 24\% & 86\% & 52\% (62\%) \\
Instructions & 5 & 73\% & 87\% & 80\% \\
Reading & 4 & 25\% (100\%) & 100\% & 100\% \\
\bottomrule
\end{tabularx}
\end{table}

\begin{table}[H]
\centering\small
\begin{tabular}{@{}lrr@{}}
\toprule
Image category & Items & Vision model \\
\midrule
Colour, shape, position, size & 13 & 100\% \\
Bar chart & 2 & 83\% \\
Counting & 5 & 73\% \\
Two images & 2 & 50\% \\
OCR (pixel font) & 5 & 20\% \\
\bottomrule
\end{tabular}
\end{table}

\begin{table}[H]
\centering\small
\begin{tabularx}{\textwidth}{@{}YllR@{}}
\toprule
Condition & Tokens, median / p90 / max & Time, median / p90 / max & Cut off \\
\midrule
Chat preset & 3 / 17 / 158 & 0.06 / 0.25 / 1.99\,s & 0 \\
Reasoning preset & 156 / 266 / 4096 & 1.92 / 3.15 / 62.95\,s & 1 \\
Vision, image items & 2 / 6 / 7 & 0.06 / 0.40 / 0.77\,s & 0 \\
Vision, text items & 3 / 6 / 600 & 0.05 / 0.09 / 6.54\,s & 0 \\
\bottomrule
\end{tabularx}
\end{table}

Image prompts were a median of 1090 tokens, almost all of it the image. Median decode
speed reported by the server was 69--90 tokens/s. One server launch per model was
recorded: 3.4\,s for the text model and 4.1\,s for the vision model.

\clearpage
\subsection*{Typical failures}
Raw replies, all three seeds unless noted.

\begin{table}[H]
\centering\small
\begin{tabularx}{\textwidth}{@{}llY@{}}
\toprule
Item (condition) & Key & Reply \\
\midrule
Capital of Australia (chat) & Canberra & \code{Sydney} \\
\code{len("oakgpt")} (chat) & 6 & \code{12} \\
\code{"ab" * 3} (chat) & ababab & \code{abcabcabc} \\
3 * (4 + 5) $-$ 6 (chat) & 21 & \code{14} \\
Sweden's largest lake (chat) & V\"anern & \code{Svanen}, \code{Stockholmsviken},
  \code{Baltic Sea} (one per seed) \\
``apple'' in Swedish (all text conditions) & \"apple & \code{appel} \\
14:20 + 95 minutes (chat) & 15:55 & \code{14:20 + 95 minutes = 15:15} \\
\code{print(7 // 2, 7 \% 2)} (reasoning) & 3 1 & \code{3,1} (format) \\
Reverse ``stone'' (reasoning, 2 seeds) & enots & \code{E-N-O-T-S} \\
Code \code{K7Q-29X} (vision) & K7Q-29X & \code{0728}, \code{0x789a} \\
\code{MAX SPEED 50} (vision) & 50 & \code{30} \\
More squares, first or second image (vision) & second & \code{first} \\
\bottomrule
\end{tabularx}
\end{table}

\section{Interpretation}\label{sec:interpretation}

\textbf{The text model's chat mode is weak.} Without thinking, Qwen3 1.7B gets half of
these simple questions right. Some errors are confident and plain wrong (Sydney,
\code{12}), which matches the C1 observation of hallucination. Its Swedish is poor: 1 of
18 Swedish replies was strictly correct. With thinking, the same model reaches 92\%,
with only a few format slips. This is the clearest result in the report, and its interval
is far from zero.

\textbf{Product consequence, not yet applied.} The UI defaults to the chat preset, which
is fast but often wrong on these items. The reasoning preset is about 30 times slower per
short answer, but still under 2\,s at the median, and much more accurate here. Whether to
switch the default is a decision for the project owner. Long answers would take
proportionally longer, and a 63\,s worst case was observed.

\textbf{The vision model is a reasonable text model.} Answering text directly, it beats
the text model's chat mode, most visibly on Swedish, code and reading. That supports
keeping image conversations on the vision model. It also raises the question of whether
the vision model should answer direct (non-thinking) text questions too. This set is too
small to decide that, but it can be tested on \code{chat\_eval\_v2}.

\textbf{Image perception is uneven.} Simple visual properties were always right.
Counting fails near 7 objects, and one of two two-image items failed on every seed. OCR
results need care. The same model read ordinary UI text verbatim in C1.1, and here it
read \code{OPEN 24 HOURS} and the chart letters. It failed on the codes and numbers in
the 5$\times$7 font. So the 20\% most likely reflects this unusual blocky font plus small
digits rather than OCR in general. A v2 set should add text rendered in a normal
typeface, without changing v1.

\section{Limitations}
\begin{itemize}
\item \textbf{Small set.} 67 items. Per-category numbers rest on 2--7 items, and overall
  intervals are 25--30 points wide in chat mode.
\item \textbf{One author.} Templated items are not independent, and question clusters
  make the intervals somewhat optimistic~\cite{miller}.
\item \textbf{One run on one host.} Seeds fix sampling, but GPU arithmetic is not
  guaranteed to be bit-identical across runs.
\item \textbf{Synthetic images.} Shapes and pixel text test basic perception, not
  photographs, handwriting or documents.
\item \textbf{Contamination.} The pretrained models may have seen similar facts and
  arithmetic. The set measures behaviour, not novelty.
\item \textbf{Scoring heuristics.} Strict scoring penalises answers that are right but in
  an unexpected form. Lenient scoring can pass hedged replies. Both are recorded, along
  with the raw text.
\item \textbf{Not a benchmark.} This is a regression and comparison set for this project,
  not a standardised benchmark. No comparison to published scores is valid.
\end{itemize}

\section{Reproduction}
\begin{Verbatim}[fontsize=\footnotesize,frame=single,framesep=3mm]
.\.venv\Scripts\python scripts/build_chat_eval_v1.py <empty-dir>   # identical files
.\.venv\Scripts\oakgpt chat-eval --config configs/chat/qwen3_q4km.json `
  --dataset data/eval/chat_eval_v1/items.jsonl --output runs/chat/eval/rerun.json
\end{Verbatim}
The full run needed 5.9 minutes of generation, 5.1 of them in the reasoning condition,
plus model loading and file hashing.

\begin{thebibliography}{6}

\bibitem{biderman} S. Biderman et al., ``Lessons from the Trenches on Reproducible
  Evaluation of Language Models,'' arXiv:2405.14782, 2024.
  \url{https://arxiv.org/abs/2405.14782}

\bibitem{squad} P. Rajpurkar, J. Zhang, K. Lopyrev, P. Liang, ``SQuAD: 100,000+
  Questions for Machine Comprehension of Text,'' EMNLP 2016, arXiv:1606.05250.
  \url{https://arxiv.org/abs/1606.05250}

\bibitem{judge} L. Zheng et al., ``Judging LLM-as-a-Judge with MT-Bench and Chatbot
  Arena,'' NeurIPS 2023 Datasets and Benchmarks, arXiv:2306.05685.
  \url{https://arxiv.org/abs/2306.05685}

\bibitem{qwen3card} Qwen Team, ``Qwen3-1.7B'' model card, section ``Best Practices.''
  \url{https://huggingface.co/Qwen/Qwen3-1.7B}

\bibitem{efron} B. Efron, R. J. Tibshirani, \emph{An Introduction to the Bootstrap},
  Chapman \& Hall, 1993.

\bibitem{miller} E. Miller, ``Adding Error Bars to Evals: A Statistical Approach to
  Language Model Evaluations,'' arXiv:2411.00640, 2024.
  \url{https://arxiv.org/abs/2411.00640}

\end{thebibliography}

\end{document}
